\section{Open Questions}
\label{sec:openq}

The following five questions are genuinely open after this replication ---
answering any of them would materially advance the model beyond what the paper
demonstrates. Each is grounded in a specific gap between what the paper claims
and what our replication could actually test.

\subsection*{Q1. Single-cell heterogeneity vs. population-mean synchrony}
\textbf{Question.} Does the mean-field ODE core reproduce single-cell p53
pulse heterogeneity (Geva-Zatorsky--style bimodal, asynchronous pulsing) once
the Haseltine--Rawlings stochastic layer is added, or does the model, even in
its full hybrid form, still produce over-synchronized population traces
relative to real MCF-7 / U-2 OS live-cell imaging?

\textbf{Basis.} The paper's Figs.~2--3 report mean$\pm$IQR over 1000 cells but
do not compare pulse-timing variance or asynchrony coefficient against
Geva-Zatorsky (Mol.\ Syst.\ Biol.\ 2006/2010) single-cell data. Our
replication is mean-field only and cannot answer this. The shared-parameter
architecture (all 1000 cells identical rate constants) is a known source of
artificial synchrony.

\textbf{Next steps.} Implement MOESM2 Haseltine--Rawlings scheme; run 1000-cell
simulations; compute (a) CV of first-pulse peak time, (b) auto-correlation
decay of $p53_n(t)$, (c) fraction with $\geq$2 vs 1 vs 0 pulses in 24\,h.
Compare against Geva-Zatorsky and Purvis et al.\ (2012). If synchrony too high,
add per-cell log-normal parameter jitter ($\sigma \sim 0.2$ on Wip1/Mdm2
production).

\subsection*{Q2. Extension to ATR/CHK1 (S-phase / replication stress)}
\textbf{Question.} How does the model behave under S-phase-specific stress
where ATR/CHK1 (not ATM/CHK2) is the dominant DDR kinase --- and is the current
ATM-only architecture even applicable to replication-stress agents (HU,
aphidicolin, gemcitabine) that clinicians combine with radiotherapy?

\textbf{Basis.} MOESM1 Table~1 has no ATR/CHK1 species. Real DDR involves
ATM--ATR crosstalk (Cimprich \& Cortez 2008; Blackford \& Jackson 2017) that
modulates p53 in S phase. This restricts the model to G1/G0 DSB-only scenarios
--- a scope restriction the paper does not flag.

\textbf{Next steps.} Extend with a minimal ATR/CHK1 module (5--6 states: ATR,
ATRp, CHK1, CHK1p, CDC25A, replication-fork stress input $R_s$).
Parameterize from Lyu et al.\ (2019, iScience) or Wortel et al.\ (2017).
Test whether extended model reproduces (i) p53 stabilization under HU alone,
(ii) HU$+$IR synergy in S-phase-synchronized cells. Compare against
Powell/Iliakis clonogenic data.

\subsection*{Q3. Parameter identifiability of the mechanistic conclusions}
\textbf{Question.} Are the fitted parameter values quantitatively
identifiable, or one point of a large equivalence class producing
indistinguishable p53 and NF-$\kappa$B trajectories --- and if the latter,
which of the paper's mechanistic conclusions (Wip1 as dominant DDR switch-off;
TNF$\alpha$ radio-protection via NF-$\kappa$B) survive an identifiability
analysis?

\textbf{Basis.} $\sim$110 rate constants against a handful of Western-blot
time courses plus one clonogenic curve --- textbook non-identifiability
(Raue et al.\ 2009; Villaverde 2019). MOESM8 gives local sensitivity plots,
not profile likelihoods or Fisher information. Our replication reproduces the
trajectories but cannot claim the mechanism.

\textbf{Next steps.} Profile-likelihood analysis on the deterministic core
(dMod / PEtab pipeline) fitting only the western-blot time-course data.
Compute CIs on four signature params: $k_{\mathrm{wip1}}$,
$k_{\mathrm{mdm2}}$, $k_{\mathrm{deg,p53}}$, $k_{\mathrm{ikk}}$.
Mechanism-preserving reparameterization test: does raising $k_{\mathrm{mdm2}}$
while lowering $k_{\mathrm{wip1}}$ reproduce Wip1-RNAi phenotype? If yes, the
``Wip1-as-dominant-switch-off'' claim is not identifiable.

\subsection*{Q4. Low-dose stochastic activation regime}
\textbf{Question.} How does the model behave at low, physiologically-relevant
IR doses (0.05--0.5\,Gy, single-DSB regime) where DSB formation is stochastic
Poisson and cell-fate is dominated by rare-event dynamics --- does the
mean-field core meaningfully underestimate low-dose hypersensitivity
(LDHS / HRS)?

\textbf{Basis.} Dose sweep starts at 2\,Gy; DSB count is a continuous
quantity (mean-field 2.3 DSBs @ 4\,Gy vs 24 stochastic). At $\leq$0.5\,Gy the
discrete-vs-continuous gap is qualitative. Published LDHS (Joiner, Marples,
Lambin 2001) is not addressed. Our claim~\#13 gap is the warning sign.

\textbf{Next steps.} Extend Haseltine--Rawlings hybrid to Poisson-sample DSB
creation at 0.05, 0.1, 0.2, 0.3, 0.5, 1.0\,Gy. Per-cell Kracikova cell-fate;
plot survival. Compare (a) shape (LQ vs induced-repair) against Marples \&
Joiner 1993 V79 data, (b) pulse-count bimodality (Bakkenist \& Kastan 2003
single-DSB prediction). If mean-field misses LDHS, argue stochastic-first at
low dose.

\subsection*{Q5. Tumor cell-line-specific mutation contexts}
\textbf{Question.} How does the model's cell-fate landscape shift when key
p53-pathway genes are mutated in a tumor-cell-line-specific way (TP53 mutant,
MDM2-amplified, ATM-null, constitutive NF-$\kappa$B) --- does it correctly
rank the radio-sensitivity of A549 (wt-p53), H1299 (p53-null), MCF-7
(wt-p53), HCT116 (p53-KO isogenic pair), U-87 (wt-ATM vs ATM-KD)?

\textbf{Basis.} Paper validates only on Ctr-RNAi vs Wip1-RNAi in one
MCF-7-like background. It never tests whether the framework generalizes to
p53-mutant or ATM-null tumor backgrounds that dominate real oncology practice.
Arguably the single largest gap between the paper's mechanistic value and
clinical relevance.

\textbf{Next steps.} For each of $\{$TP53-null, MDM2-amp ($5\times k_{\mathrm{mdm2}}$),
ATM-null ($k_{\mathrm{ATM,act}} \to 0$), constitutive NF-$\kappa$B ($5\times$
baseline)$\}$, run 6-dose sweep; predict clonogenic ranking. Compare against
HCT116 $p53^{+/+}$ vs $p53^{-/-}$ (Bunz et al.\ 1999), A549 vs H1299, MCF-7
vs T47D, U-87 wt vs U-87 ATM-KD. Score via Spearman rank correlation. If
ranking fails, missing ingredients are almost certainly (i) DNA-repair-pathway
choice (HR vs NHEJ) which is not modelled, and (ii) mitotic-catastrophe fate
which is lumped into apoptosis here.
